\documentclass[a4paper,10pt]{article}
\usepackage{lamia-pedagogy}

\begin{document}
% PAGE 1 : Compétence 01 - Compréhension et communication orales
% =========================================================================
\ficheHeader{1 Heure}{\makebox[2.5cm]{\dotfill}}{1\textsuperscript{ère} AM}{2026 / 2027}

\vspace{-3pt}
\competenceBanner{\faHeadphones\ Compétence 01 : Compréhension et communication orales}

\begin{minipage}[t]{0.48\textwidth}
\begin{tcolorbox}[
  enhanced,
  colback=secondarylight,
  colframe=secondary,
  arc=1.5mm,leftrule=3pt,rightrule=0.5pt,toprule=0.5pt,bottomrule=0.5pt,
  top=3pt,bottom=3pt,left=6pt,right=6pt,
  title={\small\textbf{\faBullseye\ Objectifs d'apprentissage}},
  coltitle=secondary!90!black,attach boxed title to top left={yshift=-2mm,xshift=4mm},
  boxed title style={colback=white,boxrule=0.5pt,colframe=secondary,arc=1mm,top=1pt,bottom=1pt}
]
\vspace{1pt}
{\small
\begin{itemize}[leftmargin=*,noitemsep,topsep=1pt]
  \item[\faCheckCircle] Identifier le thème de la situation de communication.
  \item[\faCheckCircle] Identifier les unités de sens.
  \item[\faCheckCircle] S'exprimer en fonction de la situation de communication.
\end{itemize}
}
\end{tcolorbox}
\end{minipage}
\hfill
\begin{minipage}[t]{0.48\textwidth}
\begin{tcolorbox}[
  enhanced,
  colback=accentlight,
  colframe=accent,
  arc=1.5mm,leftrule=3pt,rightrule=0.5pt,toprule=0.5pt,bottomrule=0.5pt,
  top=3pt,bottom=3pt,left=6pt,right=6pt,
  title={\small\textbf{\faExclamationTriangle\ Difficultés constatées}},
  coltitle=accent!90!black,attach boxed title to top left={yshift=-2mm,xshift=4mm},
  boxed title style={colback=white,boxrule=0.5pt,colframe=accent,arc=1mm,top=1pt,bottom=1pt}
]
\vspace{1pt}
{\small
\begin{itemize}[leftmargin=*,noitemsep,topsep=1pt]
  \item[\faTimesCircle] Difficulté à repérer le thème lors de l'écoute du support.
  \item[\faTimesCircle] Confusion entre thème et informations secondaires.
\end{itemize}
}
\end{tcolorbox}
\end{minipage}

\vspace{1pt}
\begin{tcolorbox}[
  enhanced,
  colback=white,
  colframe=cardborder,
  arc=1.5mm,
  top=2pt,bottom=2pt,left=6pt,right=6pt,
  title={\scriptsize\textbf{\faUsers\ Élèves concernés par la remédiation}},
  coltitle=primary,attach boxed title to top left={yshift=-1.8mm,xshift=3mm},
  boxed title style={colback=primarylight,boxrule=0.4pt,colframe=cardborder,arc=1mm,top=0.5pt,bottom=0.5pt}
]
\vspace{1pt}
{\scriptsize
\begin{tabularx}{\linewidth}{@{}l X l X@{}}
1. & \makebox[\linewidth]{\dotfill} & 6. & \makebox[\linewidth]{\dotfill} \\
2. & \makebox[\linewidth]{\dotfill} & 7. & \makebox[\linewidth]{\dotfill} \\
3. & \makebox[\linewidth]{\dotfill} & 8. & \makebox[\linewidth]{\dotfill} \\
4. & \makebox[\linewidth]{\dotfill} & 9. & \makebox[\linewidth]{\dotfill} \\
5. & \makebox[\linewidth]{\dotfill} & 10. & \makebox[\linewidth]{\dotfill} \\
\end{tabularx}
}
\end{tcolorbox}

\vspace{1pt}
\begin{tcolorbox}[
  enhanced,
  colback=secondary!5,
  colframe=secondary!40,
  arc=1.5mm,
  top=3pt,bottom=3pt,left=7pt,right=7pt
]
\textbf{\faCompass\ Déroulement de l'activité (Mise en situation) :}
\begin{itemize}[leftmargin=*,noitemsep,topsep=1pt]
  \item[\faCommentDots] Avez-vous déjà visité un zoo avec vos parents ?
  \item[\faCommentDots] Si oui, dites ce qu'on trouve dedans ?
\end{itemize}
\end{tcolorbox}

\vspace{1pt}
\begin{tcolorbox}[
  enhanced,
  colback=amber!5!white,
  colframe=primary!70,
  arc=2mm,
  top=3pt,bottom=3pt,left=8pt,right=8pt,
  title={\small\textbf{\faVolumeUp\ Support audio transcrit : \guillemotleft\ La gazelle \guillemotright}},
  coltitle=primary,attach boxed title to top left={yshift=-2mm,xshift=4mm},
  boxed title style={colback=white,boxrule=0.5pt,colframe=primary!70,arc=1mm,top=1pt,bottom=1pt}
]
\vspace{1pt}
{\small\itshape
La gazelle est un mammifère sauvage qui vit dans le désert. Elle est herbivore, elle se nourrit d'herbes, de plantes et de feuilles d'arbres. Son poil est de couleur rousse sur le dos et blanche sur le ventre. Elle a une queue noire et quatre pattes fines terminées par des sabots. Sa petite tête porte deux cornes. La gazelle pèse jusqu'à 30 kilos.
\begin{center}
\textbf{La gazelle est un animal très rapide.}
\end{center}
}
\end{tcolorbox}

\vspace{3pt}
{\small
\setlength{\tabcolsep}{2.5pt}

% 1ère écoute
\noindent\badge[primary]{\faStepForward\ 1\textsuperscript{ère} écoute (Pré-écoute)}

\vspace{1pt}
\noindent
\begin{tabular}{@{}p{5.7cm} p{4.3cm} p{4.1cm} p{3.8cm}@{}}
$\bullet$ \textbf{Qui parle dans cette vidéo ?} & \qcmCircle a) un adulte & \qcmCircle b) un enfant & \\
$\bullet$ \textbf{De quoi parle-t-il ?} & \qcmCircle a) du lion & \qcmCircle b) de la gazelle & \qcmCircle c) du fennec \\
$\bullet$ \textbf{La gazelle est un animal :} & \qcmCircle a) sauvage & \qcmCircle b) domestique & \\
\end{tabular}

\vspace{3pt}
% 2ème écoute
\noindent\badge[primary]{\faPlay\ 2\textsuperscript{ème} écoute (Écoute active)}

\vspace{1pt}
\noindent
\begin{tabular}{@{}p{5.7cm} p{4.3cm} p{4.1cm} p{3.8cm}@{}}
$\bullet$ \textbf{La gazelle est herbivore, elle se nourrit :} & \qcmCircle a) d'herbe & \qcmCircle b) de viande & \\
$\bullet$ \textbf{Où vit-elle ?} & \qcmCircle a) dans la savane africaine & \qcmCircle b) dans le jardin & \\
$\bullet$ \textbf{Combien pèse la gazelle ?} & \multicolumn{3}{l@{}}{\textbf{Réponse :} \makebox[11.5cm]{\dotfill}} \\
\end{tabular}

\vspace{3pt}
% 3ème écoute
\noindent\badge[primary]{\faCheckDouble\ 3\textsuperscript{ème} écoute (Post-écoute)}

\vspace{1pt}
\noindent
\begin{tabular}{@{}p{5.7cm} p{4.3cm} p{4.1cm} p{3.8cm}@{}}
$\bullet$ \textbf{Les pattes de la gazelle sont-elles :} & \qcmCircle a) fines & \qcmCircle b) grosses & \\
$\bullet$ \textbf{Sa queue est :} & \qcmCircle a) noire & \qcmCircle b) jaune & \\
\multicolumn{4}{@{}l@{}}{$\bullet$ \textbf{Comment est décrite la gazelle à la fin de la vidéo ?} \quad \textbf{Réponse :} \makebox[8.5cm]{\dotfill}} \\
\end{tabular}
}

\vspace{4pt}
\begin{tcolorbox}[
  enhanced,
  colback=secondarylight,
  colframe=secondary,
  arc=2mm,
  leftrule=3.5pt,rightrule=0.6pt,toprule=0.6pt,bottomrule=0.6pt,
  top=6pt,bottom=6pt,left=8pt,right=8pt
]
\begin{minipage}[c]{0.58\textwidth}
\textbf{\faPenSquare\ Synthèse : Complète l'énoncé suivant}\\[4pt]
{\small
\noindent La gazelle est un animal \dotfill ,\\[2pt]
elle est \dotfill ,\\[2pt]
elle pèse environ \dotfill ,\\[2pt]
elle vit dans \dotfill .
}
\end{minipage}
\hfill
\begin{minipage}[c]{0.38\textwidth}
\centering
\includegraphics[width=\linewidth,height=3.8cm,keepaspectratio]{images/gazelle_p1.png}
\end{minipage}
\end{tcolorbox}

\newpage
% =========================================================================
% PAGE 2 : Fiche élève à découper (2 colonnes A5)
% =========================================================================

\begin{tikzpicture}[remember picture, overlay]
\draw[dashed,thick,gray!60] ([yshift=-1cm]current page.north) -- ([yshift=1cm]current page.south);
\node[fill=white,rotate=90,inner sep=3pt] at (current page.center) {\textcolor{gray!70}{\faCut\ \scriptsize\textit{Ligne de découpe centrale}}};
\end{tikzpicture}

\newcommand{\modernGazelleCard}{%
\begin{minipage}[t][27.6cm][t]{0.96\linewidth}
\begin{tcolorbox}[
  enhanced,
  colback=primarylight,
  colframe=primary,
  arc=2mm,
  top=3pt,bottom=3pt,left=6pt,right=6pt
]
\textbf{\faFile*\ Fiche Élève : La gazelle} \hfill {\scriptsize\badge[primary]{1\textsuperscript{ère} AM}}\\[2pt]
{\scriptsize
\faUsers\ \textbf{Nom du groupe :} \makebox[4.5cm]{\dotfill}\\[2.5pt]
\textbf{Membres du groupe (4 élèves) :}\\[1.5pt]
\begin{tabularx}{\linewidth}{@{}l X l X@{}}
\textbf{1.} & \makebox[\linewidth]{\dotfill} & \textbf{2.} & \makebox[\linewidth]{\dotfill} \\
\textbf{3.} & \makebox[\linewidth]{\dotfill} & \textbf{4.} & \makebox[\linewidth]{\dotfill} \\
\end{tabularx}
}
\end{tcolorbox}

\vspace{1pt}
\begin{center}
\badge[primary]{\faQuestionCircle\ Questionnaire de compréhension orale}
\end{center}

\vspace{1pt}
{\scriptsize
\textbf{\faStepForward\ 1\textsuperscript{ère} écoute :}\\
\textbf{1/} Qui parle dans cette vidéo ?
\cardOptTwo{un adulte}{un enfant}

\textbf{2/} De quoi parle-t-il ?
\cardOptThree{du lion}{de la gazelle}{du fennec}

\textbf{3/} La gazelle est un animal :
\cardOptTwo{sauvage}{domestique}

\vspace{2pt}
\textbf{\faPlay\ 2\textsuperscript{ème} écoute :}\\
\textbf{1/} La gazelle est herbivore, elle se nourrit :
\cardOptTwo{d'herbe}{de la viande}

\textbf{2/} Où vit la gazelle ?
\cardOptTwo{dans la savane}{dans le jardin}

\textbf{3/} Combien pèse la gazelle ?\par\vspace{1pt}
\noindent\makebox[\linewidth][l]{\hspace*{2mm}\textbf{Rép. :} \dotfill}\par\vspace{2pt}

\vspace{2pt}
\textbf{\faCheckDouble\ 3\textsuperscript{ème} écoute :}\\
\textbf{1/} Les pattes de la gazelle sont-elles :
\cardOptTwo{fines}{grosses}

\textbf{2/} Sa queue est :
\cardOptTwo{noire}{jaune}

\textbf{3/} Comment est décrite la gazelle à la fin de la vidéo ?\par\vspace{1pt}
\noindent\makebox[\linewidth][l]{\hspace*{2mm}\textbf{Rép. :} \dotfill}\par\vspace{2pt}
}

\vspace{2pt}
\begin{tcolorbox}[
  enhanced,
  colback=secondarylight,
  colframe=secondary,
  arc=1.5mm,
  top=2.5pt,bottom=2.5pt,left=5pt,right=5pt,
  title={\scriptsize\textbf{\faPen\ Synthèse : Complète l'énoncé suivant}},
  coltitle=secondary!90!black,attach boxed title to top left={yshift=-1.8mm,xshift=3mm},
  boxed title style={colback=white,boxrule=0.4pt,colframe=secondary,arc=1mm,top=0.5pt,bottom=0.5pt}
]
\vspace{1pt}
{\scriptsize
\noindent La gazelle est un animal \dotfill ,\\[2pt]
elle est \dotfill ,\\[2pt]
elle se nourrit de \dotfill ,\\[2pt]
elle vit dans la \dotfill .
}
\end{tcolorbox}

\vspace{1pt}
\begin{center}
\includegraphics[width=0.98\linewidth,height=7.2cm,keepaspectratio]{images/gazelle_schema.png}
\end{center}
\end{minipage}%
}

\noindent
\begin{minipage}[t]{0.485\textwidth}
\modernGazelleCard
\end{minipage}
\hfill
\begin{minipage}[t]{0.485\textwidth}
\modernGazelleCard
\end{minipage}

\newpage

\end{document}
